\documentclass[sigconf,nonacm]{acmart}

\usepackage{booktabs}
\usepackage{multirow}

\settopmatter{printacmref=false}
\renewcommand\footnotetextcopyrightpermission[1]{}
\pagestyle{plain}

\newcommand{\fb}{\textsc{frostbit}}
\newcommand{\us}{\,\textmu s}
\newcommand{\ns}{\,ns}

\begin{document}

\title{\fb: Compiling Boolean Filters over Compressed Bitmaps into
Allocation-Free, Provably-Sized Execution Plans}

\author{William Liu}
\affiliation{%
  \institution{Exa}
  \city{San Francisco}
  \country{USA}}
\email{wliu@exa.ai}

\begin{abstract}
Compressed bitmaps in the Roaring format are the de facto standard for
boolean filtering in search and analytics engines, and a decade of work has
produced highly tuned SIMD kernels for intersecting, uniting, and subtracting
individual containers. We observe that the remaining performance is not in
the kernels: it is in everything an engine does \emph{around} them ---
allocating and growing result buffers, re-deriving representation decisions
per operation, streaming containers that cannot survive the query, and
choosing fold orders blind to the memory system.

\fb{} is a bitmap engine built on a single organizing principle: a boolean
filter expression is \emph{compiled, once, into a complete execution
manifest} by an analysis pass that pre-computes every decision the fold
would otherwise make --- the exact set of output containers per operator, a
proven byte ceiling for each (so execution contains no allocation or growth
path at all), each key's accumulator representation, the fold order, and a
container-granular liveness mask that skips every 64\,KiB block that cannot
survive the root of the expression. The executor is deliberately trivial: it
streams the provably minimal set of containers through pre-sized,
thread-pooled buffers, and its writes \emph{are} the output --- results
serialize in place, so a steady-state query performs zero heap allocations.

On an end-to-end corpus of 25{,}000 randomly-shaped filter trees (up to 100
leaves and 15 levels), \fb{} evaluates the full corpus $3.4\times$ faster
than the standard Roaring library and $2.25\times$ faster than its
nightly-only SIMD build, while winning every operation shape in a 2--16-way
microbenchmark sweep by $1.2$--$15\times$; container-granular pruning alone
accelerates selective filters by up to $33\times$ and short-circuit
evaluation resolves contradictory filters $2{,}900\times$ faster. \fb{} is
the production filtering substrate for all boolean filtering at Exa, a
web-scale search engine.
\end{abstract}

\maketitle

\section{Introduction}

Every serving system we know of that filters documents at scale eventually
converges on the same data structure: a compressed bitmap in the Roaring
format~\cite{roaring2016}, partitioned into per-64Ki-value containers that
adaptively choose between sorted arrays, packed bitmaps, and run-length
encodings. The container \emph{kernels} --- the tight loops that intersect
two sorted arrays or AND two 1024-word bitmaps --- have been optimized to
the point of diminishing returns by a long line of SIMD work
\cite{schlegel2011,inoue2014,lemire2016simd,han2018,fesia2020}.

Our starting observation, from operating a web-scale search engine, is that
production filter latency is not dominated by those kernels. It is dominated
by the scaffolding: the allocator invoked for every intermediate result; the
representation decisions (array or bitmap? run or not?) re-derived inside
every operation; containers faithfully streamed, decoded, and merged even
when a cheap argument shows they cannot contribute to the final result; and
fold loops whose memory access patterns defeat the prefetcher at exactly the
fan-ins production queries use.

\fb{} attacks the scaffolding. The design premise is that a filter
expression --- an arbitrary tree of AND, OR, and DIFF over frozen bitmaps
--- should be \emph{analyzed once, completely}, and that the analysis should
be strong enough that execution becomes a straight-line replay with no
decisions left to make. Concretely, the analysis pass pre-computes:

\begin{itemize}
\item \textbf{the exact output container set of every operator} in the tree,
  with a \emph{proven} byte ceiling for each output slot. The executor
  therefore contains no reallocation or growth path --- not as an
  optimization, but as an invariant the planner discharges. We refer to
  these ceilings as \emph{capacity clamps} (\S\ref{sec:analysis});
\item \textbf{every representation decision}: which keys accumulate as
  arrays, bitmaps, or native run containers, including a fan-in-aware
  promotion rule derived from a calibrated cost model rather than a fixed
  cardinality threshold (\S\ref{sec:dispatch});
\item \textbf{the fold order of each operator} --- key-major or
  partner-major --- selected by comparing the plan's proven accumulator
  footprint against the cache, because the two orders are memory-optimal in
  disjoint regimes (\S\ref{sec:foldorder});
\item \textbf{a container-granular liveness mask} (``hole-punching''): an
  upward pass computes each subtree's surviving key set, and every cursor in
  the plan skips dead 64\,KiB blocks before a single byte of them is
  streamed. Together with subtree short-circuit guards compiled into the
  plan, the executor streams the provably minimal container set at the same
  granularity Roaring itself uses (\S\ref{sec:shortcircuit});
\item \textbf{all working memory}: every buffer the fold will touch is sized
  from the plan and drawn from per-thread pools; results serialize in place
  inside the working arena, so the fold's writes are the output's bytes and
  a steady-state query never calls the allocator (\S\ref{sec:memory}).
\end{itemize}

None of the individual kernels in \fb{} is novel --- deliberately so; we
verified against a systematic survey of the literature (310 collected
papers) that our merge kernels match the known state of the art, and we show
where measurement told us to deviate from it. The contribution is the layer
above: to our knowledge no published system compiles boolean container-tree
queries into allocation-free, capacity-proven, liveness-pruned execution
manifests, and no published system selects fold order from a static memory
footprint proof. The result is an engine that starts from parity with the
best-tuned kernels and then wins everywhere the scaffolding used to be.

This paper makes the following contributions:
\begin{enumerate}
\item the plan-as-manifest architecture: single-pass, construction-fused
  analysis producing complete execution manifests with proven capacity
  clamps, reusable across queries (\S\ref{sec:analysis});
\item cost-model dispatch at three levels --- kernel tier, representation
  promotion, and fold order --- each calibrated and ablated on measured
  constants, including a fan-in-aware union promotion rule and, we believe,
  the first cache-footprint-driven fold-order selection (\S\ref{sec:dispatch},
  \S\ref{sec:foldorder});
\item automatic container-granular pruning and compiled short-circuit
  guards, which reduce streamed bytes to the provable minimum at container
  granularity (\S\ref{sec:shortcircuit});
\item an allocation-free execution discipline, including in-place
  serialization where the arena buffer becomes the result
  (\S\ref{sec:memory});
\item an evaluation against both distributed configurations of the standard
  Rust Roaring library, including a 25{,}000-tree corpus, complete
  per-technique ablations with negative results, and a per-tree audit
  methodology (\S\ref{sec:eval});
\item deployment experience as the filtering substrate of a production
  web-scale search engine (\S\ref{sec:deployment}).
\end{enumerate}

\section{Background and Related Work}
\label{sec:related}

\textbf{Compressed bitmaps.} Word-aligned RLE schemes (BBC, WAH and its
descendants PLWAH, COMPAX, VAL-WAH)~\cite{wah2006} trade decode cost for
space; Roaring~\cite{roaring2016,roaring2018} replaced them in most modern
systems with a two-level design --- a u16-keyed directory of per-64Ki
containers stored as sorted arrays ($\le$4096 values), 8\,KiB bitmaps, or
run lists --- that supports random access and SIMD kernels. \fb{} adopts
Roaring's container model wholesale (and is wire-convertible to it); our
frozen format adds a structure-of-arrays directory and 64-byte container
alignment so bitmaps can be operated on directly from an \texttt{mmap}.

\textbf{Set intersection kernels.} Sorted-array intersection spans the
adaptive lineage~\cite{demaine2000} and the SIMD lineage: shuffle-based
compare-all-pairs~\cite{schlegel2011,lemire2016simd}, branch-reduced
merging~\cite{inoue2014}, byte-filtered variants~\cite{han2018}, and
segment-bitmap prefilters~\cite{fesia2020}. Cost-model algorithm selection
appears per-query in web-search list intersection~\cite{kim2019}. \fb's
kernels are the portable synthesis of this line (compare-all-pairs with
both-operand rotation, Inoue-style merge networks for union, table-driven
compaction), dispatched per container pair by explicit cost models; our
survey of this literature (\S\ref{sec:eval}) found no published kernel our
measured throughput trails.

\textbf{Bitmaps in systems.} Bitmap and signature indexes power analytics
and search engines from Model 204 and FastBit through Druid, Pinot, and
BitFunnel~\cite{bitfunnel2017}; column stores accelerate scans with
bit-parallel layouts (BitWeaving, ByteSlice). These systems document the
container structures; none, to our knowledge, documents compiling filter
\emph{expressions} into pre-proven execution plans, which is precisely the
layer \fb{} contributes. UpBit and CUBIT address updatability, which \fb{}
deliberately does not: our bitmaps are frozen artifacts of an indexing
pipeline, which is what makes whole-plan static analysis sound.

\section{System Overview}
\label{sec:overview}

A \fb{} bitmap is a single contiguous buffer: a 16-byte header, a
structure-of-arrays container directory (key, type, cardinality, offset ---
8 bytes per container), and a data section in which bitmap containers sit on
64-byte boundaries. Buffers are validated once at the trust boundary and
thereafter read zero-copy, including directly from memory maps. A builder
ingests ascending values; conversions to and from Roaring are container
transcodes.

Queries are expressions over frozen bitmaps: leaves (borrowed views or
shared owned bitmaps) combined by AND, OR, and DIFF combinators.
\emph{Constructing the expression is the analysis}: each combinator fuses
its children's plans into a flat post-order program, flattening same-operator
chains (\texttt{And(And(a,b),c)} becomes one 3-way fold), and computing
everything \S\ref{sec:analysis} describes. \texttt{materialize()} replays
the program; the plan is reused across calls, which matches the production
pattern of compiled filters evaluated against many segments.

\section{The Analysis Pass}
\label{sec:analysis}

The analysis pass is the core of the system. Its output --- the
\emph{manifest} --- makes three kinds of promises to the executor, and the
executor is correct only because all three are discharged statically.

\subsection{Shapes and capacity clamps}

For every node, bottom-up, the analyzer computes the node's \emph{shape}:
the set of container keys it can emit, and for each key an upper bound on
the output's cardinality, run count, and representation, mirroring exactly
the decision ladder the kernels use. Leaf shapes are exact (read from the
directory); interior shapes are sound over-approximations. From the shape,
each operator gets an arena plan: one slot per possible output container,
with a byte capacity \emph{proven sufficient for every state the fold can
pass through} --- the verbatim-copy case, the shrunk case, the run-split
worst case ($2 + 4\cdot(\textit{runs} + \textit{subtrahend items})$ bytes),
or a full 8\,KiB bitmap, whichever the ladder can reach. Slots are claimed
by key at execution, so a plan whose key set over-approximates reality
(as intersection plans necessarily do) remains correct with zero waste
beyond the unclaimed slots.

The consequence is an invariant, not an optimization: \emph{the execution
path contains no allocation, growth, or reallocation code at all}. The
arena asserts (in debug builds) that every recorded container fits its
clamp; an 11{,}600-case randomized stress suite exercises the proof
obligations. When we extended the system with partner-major folding
(\S\ref{sec:foldorder}), the one candidate design that would have required
loosening a clamp --- seeding intersections from the driver rather than the
per-key minimum --- was rejected for exactly that reason.

\subsection{Fan-in-aware analysis}

Because same-operator chains flatten, an operator's true fan-in is not its
child count in the source expression. The analyzer threads each spliced
child's \emph{operand weight} through shape computation, so per-key fan-in
bounds are exact upper bounds on what the executor will see. This matters
because two of the dispatch decisions below are functions of fan-in, and an
early version that counted tree children instead of flattened operands
produced plans whose kernels disagreed with their clamps --- caught
immediately by the stress suite, and fixed by making the plan's clamp the
\emph{single authority} for representation choice: the kernel never
re-derives a decision the planner already proved.

\section{Kernel and Representation Dispatch}
\label{sec:dispatch}

\fb{} mixes algorithms at every level where measurement showed a crossover,
and encodes each crossover as an explicit, commented cost model.

\textbf{Three-tier sorted-array kernels.} Balanced pairs
($|B|\le 4|A|$) take a shuffle-merge: an 8-lane compare-all-pairs using
rotations of \emph{both} operands (5 permutations instead of 7), with
branchless table-driven output compaction. Moderately skewed pairs take a
galloping broadcast-scan (one lane broadcast against a galloped window of
the larger side). Heavily skewed pairs take a scalar galloping binary
search, but only when its $|A|\log|B|$ work clears the scan's $|B|/8$ linear
advance \emph{by a margin} --- binary search mispredicts, so on-paper parity
is not enough. Every constant in these predicates was set by measurement,
and \S\ref{sec:eval} reports the experiments --- including the ones that
failed.

\textbf{Fan-in-aware union promotion.} Union accumulators outgrow arrays.
The classical rule promotes to a bitmap at a fixed cardinality threshold
(256 in CRoaring-derived engines). But the array path's cost is
$\Theta(\textit{sum}\cdot(n{+}1)/2)$ streamed elements at fan-in $n$ (each
merge re-streams the accumulator), while the bitmap path is a flat clear +
scatter + popcount. Solving with measured constants gives a promotion rule
that is a function of both: never at $n\le3$, above ${\sim}1000$ summed
values at $n{=}4$, falling as fan-in grows. The fixed threshold cannot be
made correct: tuned low it destroys low-fan-in expressions (AND of small
OR-groups --- we measured $89\to115\us$ on that shape), tuned high it
abandons wide unions (we measured $203\us$ on a 4-way union the fan-in rule
executes in $54.9\us$).

\textbf{Run containers.} Dense runs fold natively --- run$\cap$run,
run$\cup$run, and run-splitting difference operate on typed run slices
inside bitmap-sized clamps, converting to bitmaps only on proven overflow.

\section{Fold Order and the Memory System}
\label{sec:foldorder}

An $n$-way fold visits a cross product of containers: keys $\times$
operands. \emph{Key-major} order (gather all operands at a key, fold, move
to the next key) keeps one accumulator cache-hot but resumes $n$ interleaved
input streams per key. \emph{Partner-major} order (stream one operand fully,
then the next) reads every input as a long sequential stream --- the pattern
prefetchers are built for --- but revisits the whole accumulator set every
pass.

Profiling difference folds at fan-in 16 localized a $+50\%$ per-block
penalty entirely to memory: the kernel ran at $1.45\ns$/8-lane-block
standalone, $2.17\ns$ inside a 32-key fold, and $1.56\ns$ in a single-key
control --- the interleaved resumption of 16 partner streams was defeating
L1. Partner-major folding recovered it, \emph{and then dense unions showed
the mirror image}: 16 bitmap accumulators are 128\,KiB, revisiting them
every pass thrashes L1, and key-major won by 12\%.

The resolution is the system's thesis in miniature: \emph{the plan already
knows the answer}. The accumulator footprint is the sum of the proven
capacity clamps; if it is cache-resident ($\le$64\,KiB) the plan is laid out
double-buffered and the fold runs partner-major (array merges flip each slot
between two arena sides, so out-of-place kernels need no staging copies);
otherwise key-major. Intersection is the deliberate exception: a 1-key
vs.\ 32-key control shows 16-way AND scales linearly with key count
($695\ns\to21.6\us$, $31.1\times$ for $32\times$) --- after the first fold
the accumulator is a sliver and partners are galloped into, not re-streamed
--- so there is no pathology to fix, and partner-major would have required
loosening the intersection clamps. Each operator got the fold order the
measurements assigned it.

\section{Streaming the Provable Minimum}
\label{sec:shortcircuit}

The analysis pass eliminates work at three granularities, all compiled into
the manifest.

\textbf{Hole-punching (container granularity).} An upward pass computes each
node's surviving key set ($\cap$ for AND, $\cup$ for OR, left side for
DIFF). When an AND's intersected key set is provably narrower than some
child's, the planner materializes it as an 8\,KiB liveness bitset and every
leaf cursor in the plan skips dead keys in place --- dead blocks are never
streamed, decoded, or merged. The check is free (compare key counts), the
mask is derived from the plan's own shape (a stale or foreign mask is
unrepresentable), and pruning is sound because it is monotone: only keys
that cannot reach the root are dropped. A selective filter --- one narrow
4-block conjunct against 256-block disjunctions --- runs in $7.35\us$
against $241\us$ unpruned and $374\us$ for Roaring's SIMD build. This is,
to our knowledge, the first published container-granular analogue of
zone-map pruning that operates \emph{across} a boolean expression tree
rather than along a scan.

\textbf{Subtree short-circuits (operator granularity).} The compiled program
carries guard instructions: if an AND operand (or a DIFF minuend) evaluates
empty, the guard drops the partial operands and jumps past the fold ---
sibling subtrees are never evaluated. Guards use relative offsets so they
survive splicing. A contradictory filter (empty subtree beside an expensive
disjunction) resolves in $133\ns$ against $392\us$ for Roaring:
$2{,}900\times$, because the expensive branch is simply never run.

\textbf{Intra-operator short-circuits (element granularity).} Intersections
fuse population counting into the AND so an accumulator that empties stops
the key immediately; per-key folds order partners most-selective-first;
per-slot state kills dead keys across partner-major passes.

\section{Memory Discipline}
\label{sec:memory}

Every buffer in the execution path is pre-allocated and pooled per thread:
the op arena (slot storage + fixed scratch), the fold's cursor and reference
scratch, the tree evaluator's operand stack, and --- closing the loop ---
the result buffers themselves. Arenas are laid out with worst-case header
and directory room reserved at the front, so serialization \emph{compacts
payloads leftward in place} (sources never precede destinations, by the
clamp ordering) and hands the arena's own buffer out as the result; a
result-pool buffer is swapped back so the pools circulate. Tree evaluation
chains intermediate arenas directly as operator inputs --- nothing is
serialized until the root.

The consequence is measurable, not rhetorical: after warm-up, a query
performs \emph{zero} heap allocations end-to-end, verified by sampling
profiles (the standard library's allocator simply does not appear).
Roaring's evaluator, by contrast, allocates every intermediate; in our
profiles roughly 6\% of its runtime is allocator time --- which means \fb{}
starts every comparison six points ahead before any algorithmic difference
applies, and exhibits none of the allocator-induced tail variance that
motivated this design in production.

\section{Evaluation}
\label{sec:eval}

\textbf{Setup.} Criterion, 30 samples $\times$ 4\,s per benchmark
(corpus: 10 $\times$ 20\,s), Apple M-series (aarch64/NEON). Baselines: Rust
\texttt{roaring} v0.11.4 in both of its configurations --- the default build
every stable-toolchain user gets, and the nightly-only \texttt{simd} build
(its strongest form). All tree benchmarks time \fb's full pipeline including
plan construction. A two-pass runner interleaves both baselines into one
report. Correctness is enforced throughout by differential testing against
Roaring (400-tree random differentials, punched-tree differentials, run
containers, 10M-element round trips) plus the 11.6k-case clamp stress.

\subsection{Operation sweep}

Across the full 2--16-way sweep over sparse-array, dense-bitmap, and
run-container regimes, \fb{} wins \emph{every cell} against default Roaring
by $1.2$--$15\times$ (sparse-array operations by $3.6$--$15\times$: its
array kernels are scalar without the nightly feature, while \fb's SIMD
always runs). Against the nightly SIMD build, \fb{} wins everywhere except
two 16-way cells inside measurement noise ($0.98\times$, $0.99\times$).
Table~\ref{tab:ops} shows the 8-way column.

\begin{table}[t]
\caption{8-way folds, medians (\textmu s). R = Roaring default,
R\textsuperscript{simd} = nightly SIMD build.}
\label{tab:ops}
\begin{tabular}{lrrr}
\toprule
 & \fb & R & R\textsuperscript{simd} \\
\midrule
AND (sparse / dense / runs) & \textbf{17} / \textbf{12} / \textbf{2.7} & 73 / 23 / 4.6 & 25 / 22 / 4.4 \\
OR (sparse / dense / runs) & \textbf{117} / \textbf{21} / \textbf{3.0} & 1200 / 39 / 7.1 & 789 / 39 / 7.2 \\
DIFF (sparse / dense / runs) & \textbf{72} / \textbf{53} / \textbf{2.0} & 1110 / 158 / 5.0 & 80 / 155 / 5.1 \\
\bottomrule
\end{tabular}
\end{table}

\subsection{Filter-tree corpus}

The headline benchmark is a corpus of 25{,}000 deterministic random filter
trees (295{,}455 leaves; sizes log-spread 2--100 leaves, depths 2--15;
per-tree shape profiles spanning deep AND-chains, flat 100-way ORs, and
DIFF-heavy composites; 100 trees pinned to both extremes at once), evaluated
end-to-end with per-query plan construction. \fb{} clears the corpus in
\textbf{2.63\,s} (9.5K trees/s) against 9.06\,s for default Roaring
($3.4\times$) and 5.92\,s for the SIMD build ($2.25\times$). Named
production-like shapes: a conjunction of OR-groups runs $72.6\us$ vs.\
$331/102\us$; a five-way flattened conjunction $34.1\us$ vs.\ $315/80\us$; a
base-and-negation filter $17.5\us$ vs.\ $41/44\us$.

A per-tree audit (minimum-of-5 timing of every corpus tree on both engines)
is part of our methodology: at its worst point the audit surfaced 1{,}245
trees ${>}3\us$ behind the SIMD build; root-causing them produced the
hoisted bitmap-probe loop (offenders fell to 743) and ultimately the auto
hole-punching pass. We consider the audit-driven loop --- measure per-shape,
explain the worst offender structurally, fix the class --- a reusable
methodology, and we report it because each fix is an ablation with a
causal story.

\subsection{Ablations, including the failures}

Every mechanism was landed under a keep-only-if-faster discipline; the
negative results delimit the design space as sharply as the positives:

\begin{itemize}
\item \textbf{Kernel tiers}: the balanced shuffle-merge took the flagship
  DNF tree from $130\to57\us$; the register-resident rewrite of that kernel
  halved a 2-way AND ($49\to24.8\us$, ahead of Roaring's 33);
\item \textbf{fan-in promotion}: 4-way sparse OR $203\to54.9\us$ with the
  guard shapes (AND of 3-way OR-groups) \emph{unchanged} --- the fixed-256
  rule, measured, sacrifices one side or the other;
\item \textbf{fold order}: sparse 8-way DIFF $85\to72\us$
  (partner-major), dense OR unchanged only because the footprint gate
  routes it key-major --- forcing partner-major there costs 12\%;
\item \textbf{hole-punching, automatic}: whole-corpus $3.10\to2.63$\,s
  ($-15.3\%$, $p<0.05$); selective trees $33\times$;
\item \textbf{in-place serialization}: the profiling that motivated it
  showed DIFF's fold already \emph{beating} Roaring while the end-to-end
  number lost --- the entire gap was the output copy Roaring never pays;
  eliminating it flipped the cells and sped up every large-result operation
  (dense 2-way OR $7.2\to5.1\us$);
\item \textbf{branchlessness, selectively}: table-compaction beats the
  per-hit emit loop $2.2\times$ at dense overlap and is gated to cost
  nothing at sparse overlap; but a branchless scalar union merge
  \emph{regressed} $72\to91\us$ (the predictable branch beats unconditional
  work), a deferred vector-mask accumulation regressed DIFF $83\to97\us$
  (loop-carried dependency), and an adversarial-vs-predictable input sweep
  showed the merge advance branches cost \emph{zero} --- branchless
  transformations pay only where decisions are dense and data-dependent;
\item \textbf{rejected candidates}: window pre-clipping (redundant with
  galloping), eager bitmap intermediates (wins flat unions, loses
  tree-consumers), partner-major intersection (linear-scaling control shows
  nothing to win, and it would weaken the clamps).
\end{itemize}

\section{Deployment}
\label{sec:deployment}

\fb{} is the foundational bitmap structure for all boolean filtering at
Exa, a production web-scale search engine: every filter predicate ---
domain, language, date, category, and their arbitrary boolean combinations
--- compiles to a \fb{} expression evaluated against frozen per-segment
bitmaps served from memory maps. The properties this paper measures are the
properties that matter operationally: plans are built once per query and
replayed across segments; the zero-allocation hot path removes
allocator-induced tail latency; and hole-punching converts the dominant
production pattern --- one narrow conjunct against wide disjunctions ---
from the worst case into the best. The frozen (read-only) model is not a
limitation in this architecture but the enabling assumption: segments are
immutable artifacts of the indexing pipeline, so whole-plan static analysis
is sound by construction.

\section{Conclusion}

\fb{} demonstrates that after a decade of kernel optimization, the remaining
factor of $2$--$15\times$ in boolean bitmap filtering was hiding in the
scaffolding, and that the way to claim it is an analysis pass strong enough
to make execution trivial: pre-compute every output, prove every capacity,
choose every algorithm and fold order from calibrated cost models, prune
every dead container, and pre-allocate every byte --- then let a
decision-free executor stream the provable minimum. The library, its
differential test suites, its benchmark corpus, and the per-tree audit
tooling are open source.

\begin{thebibliography}{9}

\bibitem{roaring2016}
S. Chambi, D. Lemire, O. Kaser, R. Godin.
\newblock Better bitmap performance with Roaring bitmaps.
\newblock \emph{Software: Practice and Experience} 46(5), 2016.

\bibitem{roaring2018}
D. Lemire et al.
\newblock Roaring Bitmaps: Implementation of an Optimized Software Library.
\newblock \emph{Software: Practice and Experience} 48(4), 2018.

\bibitem{wah2006}
K. Wu, E. J. Otoo, A. Shoshani.
\newblock Optimizing bitmap indices with efficient compression.
\newblock \emph{ACM TODS} 31(1), 2006.

\bibitem{schlegel2011}
B. Schlegel, T. Willhalm, W. Lehner.
\newblock Fast sorted-set intersection using SIMD instructions.
\newblock \emph{ADMS@VLDB}, 2011.

\bibitem{inoue2014}
H. Inoue, M. Ohara, K. Taura.
\newblock Faster set intersection with SIMD instructions by reducing branch
  mispredictions.
\newblock \emph{PVLDB} 8(3), 2014.

\bibitem{lemire2016simd}
D. Lemire, L. Boytsov, N. Kurz.
\newblock SIMD compression and the intersection of sorted integers.
\newblock \emph{Software: Practice and Experience} 46(6), 2016.

\bibitem{han2018}
S. Han, L. Zou, J. X. Yu.
\newblock Speeding up set intersections in graph algorithms using SIMD
  instructions.
\newblock \emph{SIGMOD}, 2018.

\bibitem{fesia2020}
J. Zhang, Y. Lu, D. G. Spampinato, F. Franchetti.
\newblock FESIA: A fast and SIMD-efficient set intersection approach on
  modern CPUs.
\newblock \emph{ICDE}, 2020.

\bibitem{kim2019}
S. Kim et al.
\newblock List intersection for web search: algorithms, cost models, and
  optimizations.
\newblock \emph{PVLDB} 12(1), 2019.

\bibitem{demaine2000}
E. D. Demaine, A. L\'opez-Ortiz, J. I. Munro.
\newblock Adaptive set intersections, unions, and differences.
\newblock \emph{SODA}, 2000.

\bibitem{bitfunnel2017}
B. Goodwin et al.
\newblock BitFunnel: Revisiting signatures for search.
\newblock \emph{SIGIR}, 2017.

\end{thebibliography}

\end{document}
